\makereporttitle{2024}{ianuarie}

\makereportpreamble

\textbf{Rezumat:}

În luna ianuarie, m-am concentrat pe aprofundarea unor concepte fundamentale pentru fizica modernă, cum ar fi mecanica statistică, termodinamica și electrodinamica clasică (legile lui Maxwell), precum și teoria relativității speciale (grupul Lorentz și grupul Poincaré) și a relativității generale (geometrie diferențială și varietăți pseudo-Riemanniene).

Pe partea de cercetare, am continuat să lucrez la reproducerea și validarea unor rezultate obținute în lucrările de cercetare anterioare ale d-lui.~prof.~Băran. Folosind modulul de \href{https://docs.sympy.org/latest/modules/physics/quantum/index.html}{mecanică cuantică} din biblioteca \href{https://www.sympy.org/en/index.html}{SymPy}, am extins termenii și am cuantizat un Hamiltonian cuadrupol bosonic, provenit dintr-un model al mișcărilor nucleare, pe care l-am comparat apoi cu formula dată în „Classical and quantal chaos described by a fourth order quadrupole boson Hamiltonian”, articol publicat de V.~Băran și A.~A.~Raduta în \textit{International Journal of Modern Physics E} în 1998. Am descoperit o neconcordanță între expresia calculată simbolic de program și formula din lucrare, pentru care încă încercăm să găsim o explicație.

\vspace{0.5em}

\textbf{Cursuri și seminare științifice:}

Pentru a mă pregăti pentru evaluarea la cursul doctoral de „Fizica particulelor elementare: SM și MSSM”, ținut de C.S.~I dr.~Călin Alexa, am parcurs primele câteva capitole din cartea „Quarks and leptons: An Introductory Course in Modern Particle Physics” a lui Francis Halzen și Alan D.\ Martin, precum și o serie de video-uri și articole referitoare la evoluția modelului standard și la experimentele desfășurate în cadrul acceleratorului LHC de la CERN.

În zilele de 23 și 25 ianuarie am participat la prelegerile ținute în cadrul conferinței „Analysis and Potential”, organizată de Institutul de Matematică al Academiei Române cu ocazia zilei aniversare a lui prof.~dr.~Gheorghe Bucur. Printre altele, am urmărit prezentări despre ecuații neliniare de tip Fokker-Planck-Kolmogorov, generalizări ale formulei Trotter-Kato, ecuații diferențiale stocastice cu zgomot multiplicativ, utilizarea spinorilor în \emph{quantum computing} și ecuația de forță Lorenz.

\vspace{0.5em}

\textbf{Observațiile conducătorului științific}:

\vspace{2em}

\makereportpostamble{31 ianuarie 2024}